% TOP_PROBLEM: {"id": 2000016, "title": "Some Open Problems in Elasticity — Quasiconvexification of energy wells", "category_id": 9, "category": {"id": 9, "name": "pde", "display_name": "Partial Differential Equations", "description": "PDEs and their applications in physics and geometry.", "slug": "pde", "order_index": 9, "created_at": "2026-07-31T15:26:25.671Z"}, "status": "partially_solved"}
% FIRST_POSTED: null
% ENTRY_KIND: statement_only
% Imported from: batch08.tex; UnsolvedMath ID 2000016
% Compile twice: lualatex 2000016.tex
\documentclass[11pt]{article}
\usepackage{fontspec}
\setmainfont{Latin Modern Roman}
\usepackage{amsmath,amssymb,mathtools,unicode-math}
\setmathfont{Latin Modern Math}
\usepackage{xurl,hyperref}
\hypersetup{hidelinks,pdftitle={UnsolvedMath Problem 2000016}}
\newcommand{\sref}[2]{\hyperlink{source:#1:#2}{[S#2]}}
\newcommand{\cataloguescope}[1]{\paragraph{Mathematical question.}#1}
\begin{document}
% UnsolvedMath ID 2000016; source catalogue code AMR-019-0016
\setcounter{section}{2000015}
\section{Some Open Problems in Elasticity — Quasiconvexification of energy wells}\label{2000016}
{\small\sffamily UnsolvedMath ID 2000016 · Partial Differential Equations}
\subsection{Definitions and mathematical statement}

\paragraph{Shared notation.}$\mathbb N=\{1,2,\ldots\}$, and $\mathbb Z,\mathbb Q,\mathbb R,\mathbb C$ denote the integers, rationals, reals and complex numbers. Any other conventions are specified locally.


Write $M_+^{3\times3}=\{A\in\mathbb R^{3\times3}:\det A>0\}$ and $SO(3)=\{R:R^TR=I,\det R=1\}$. Let $P^{24}\subset SO(3)$ be the cubic rotation group, and $U(\theta)$ a positive-definite symmetric transformation strain. Let $U_1(\theta),\ldots,U_N(\theta)$ be the distinct matrices $Q^TU(\theta)Q$, $Q\in P^{24}$, with $N$ independent of temperature in the source's interval $E$ containing $\theta_c$. The normalized free energy $\psi(A,\theta)$ has zero set \[K(\theta)=\begin{cases}\displaystyle\bigcup_{i=1}^N SO(3)U_i(\theta),&\theta<\theta_c,\\\displaystyle SO(3)\cup\bigcup_{i=1}^N SO(3)U_i(\theta_c),&\theta=\theta_c.\end{cases}\] Large lattice-invariant shears, which would give infinitely many wells, are excluded in this model. A continuous function $\varphi:\mathbb R^{3\times3}\to\mathbb R$ is quasiconvex if $\varphi(A)\le\int_{(0,1)^3}\varphi(A+Dv(x))\,dx$ for every $v\in W^{1,\infty}_0((0,1)^3;\mathbb R^3)$. For compact $K$, define \[K^{qc}=\{A:\varphi(A)\le\max_{B\in K}\varphi(B)\text{ for every quasiconvex }\varphi\}.\]
\paragraph{Statement recorded in the supplied source.}
Determine $K(\theta)^{qc}$ for every $\theta\in E$ with $\theta\le\theta_c$. This includes the cubic-to-tetragonal case $U(\theta)=\operatorname{diag}(\eta_3,\eta_1,\eta_1)$, with $\eta_1,\eta_3>0$, its three permuted variants, and the addition of the austenite well $SO(3)$ at $\theta_c$. \sref{2000016}{2}
\subsection{Short English statement}
Determine the macroscopic zero-energy deformation gradients generated by the crystal energy wells, both below and at the transition temperature.
\subsection{Sources}
\begin{description}
\item[\hypertarget{source:2000016:1}{S1}] ULAM.AI and the UnsolvedMath Contributors, UnsolvedMath: A Curated Collection of Open Mathematics Problems, record 2000016 (AMR-019-0016). CC BY 4.0. \url{https://huggingface.co/datasets/ulamai/UnsolvedMath}.
\item[\hypertarget{source:2000016:2}{S2}] J. M. Ball, Some Open Problems in Elasticity (2002), Problem 16 and equation (4.1), section 4.2, pages 29–31. \url{https://people.maths.ox.ac.uk/ball/Articles%20in%20Conference%20Proceedings%20and%20Books/JMB%202002%20re%20Marsden%2060th.pdf}.
\end{description}
\label{end:2000016}
\end{document}
